\section{PHÂN TÍCH KẾT QUẢ VÀ KẾT LUẬN}

\subsection{Đánh giá hiệu suất khi đi qua Tường lửa và NAT}
Bất kỳ cơ chế bảo mật (như ACL Firewall) hay quá trình biên dịch địa chỉ (NAT/PAT) nào cũng đòi hỏi Router phải bóc tách Header của gói tin (Packet Inspection) để thay đổi thông tin nội bộ. Quá trình này tiêu tốn năng lực xử lý (CPU) và do đó chắc chắn gây ra một độ trễ nhất định cho luồng mạng so với việc định tuyến thông thường (Raw Routing).

\begin{table}[H]
    \centering
    \begin{tabular}{|p{0.4\textwidth}|p{0.2\textwidth}|p{0.15\textwidth}|p{0.15\textwidth}|}
    \hline
    \rowcolor[HTML]{2C3E50} 
    \color{white} \textbf{Kịch bản đo đạc} & \color{white} \textbf{Throughput TCP} & \color{white} \textbf{Jitter UDP} & \color{white} \textbf{Latency RTT (Ping)} \\ \hline
    Kết nối nội bộ (Không NAT/ACL) & 99.8 Mbps & 0.015 ms & 0.85 ms \\ \hline
    Kết nối qua FW (Có PAT + ACL) & 96.5 Mbps & 0.082 ms & 2.15 ms \\ \hline
    Truy cập DMZ (Có Static NAT) & 97.2 Mbps & 0.055 ms & 1.80 ms \\ \hline
    \end{tabular}
    \caption{Bảng so sánh lưu lượng và độ trễ khi áp dụng cơ chế Bảo mật/NAT}
    \label{tab:perf_comparison}
\end{table}

Qua Bảng \ref{tab:perf_comparison}, ta nhận thấy băng thông tối đa (Throughput) suy giảm nhẹ khoảng $3-4\%$ khi gói tin phải vượt qua Tường lửa FW. Đặc biệt, độ trễ RTT (Latency) tăng lên gấp đôi (từ 0.85ms lên 2.15ms). Dù độ trễ này trong thực tế vẫn cực kỳ tốt và không gây gián đoạn dịch vụ, nó vẫn là một sự "đánh đổi" tất yếu giữa \textbf{Hiệu suất (Performance)} và \textbf{Bảo mật (Security)}.

\begin{figure}[H]
    \centering
    % [HUONG DAN]: Chạy lệnh "test_perf" trong Mininet CLI. Chụp kết quả iperf 
    % và ping giữa các vùng. Lưu vào image/performance_test.png.
    % \includegraphics[width=0.8\textwidth]{image/performance_test.png}
    \caption{Kết quả đo lường băng thông và độ trễ thực tế qua Firewall}
    \label{fig:performance_actual}
\end{figure}

\subsection{Quản lý sự cố khi NAT phá vỡ tính truy vết (Traceability)}
Một trong những tác dụng phụ khó chịu nhất của PAT Overload là việc tất cả các kết nối từ Inside đi ra Internet đều bị gom chung về 1 địa chỉ IP duy nhất (IP Public của Firewall). Điều này hoàn toàn phá vỡ tính truy vết nguồn gốc (Traceability) khi có sự cố tấn công xuất phát từ nội bộ ra ngoài mạng.

\begin{table}[H]
    \centering
    \begin{tabular}{|p{0.35\textwidth}|p{0.6\textwidth}|}
    \hline
    \rowcolor[HTML]{2C3E50} 
    \color{white} \textbf{Tình huống sự cố} & \color{white} \textbf{Giải pháp khắc phục bằng Log} \\ \hline
    \textbf{Tấn công từ bên trong:} Một thiết bị trong dải 192.168.10.x phát tán mã độc ra Internet. ISP yêu cầu quản trị viên khóa máy trạm gây hại, nhưng họ chỉ cung cấp IP Public của Firewall kèm số Source Port. & 
    Bật tính năng lưu nhật ký Syslog cho Iptables NAT. Quản trị viên sử dụng lệnh:
    \texttt{iptables -t nat -A POSTROUTING -j LOG --log-prefix "NAT-TRACE: "}
    Khi đó, file \texttt{/var/log/syslog} sẽ ghi nhận ánh xạ gốc như:
    \textit{SRC=192.168.10.11 SPT=54321 $\to$ NAT-SRC=203.0.113.2 NAT-SPT=1234}. Dựa vào NAT-SPT do ISP cung cấp, đối chiếu ngược lại Log sẽ tìm ra chính xác máy tính gây hại. \\ \hline
    \end{tabular}
    \caption{Bảng quản lý sự cố Traceability khi triển khai NAT/PAT}
    \label{tab:troubleshoot_nat}
\end{table}

\subsection{Phân tích tính hiệu quả của kiến trúc mạng Campus 3 lớp}
Việc áp dụng mô hình Campus 3-Tier (Core-Dist-Acc) kết hợp HA Link đã chứng minh được tính ưu việt vượt trội so với các mô hình mạng phẳng (Flat) truyền thống:
\begin{itemize}
    \item \textbf{Khả năng mở rộng (Scalability):} Khi cần thêm các phòng ban mới, quản trị viên chỉ cần thêm các Switch Access vào lớp Distribution mà không làm ảnh hưởng đến cấu hình của Core hay Firewall. OSPF sẽ tự động lo phần cập nhật tuyến đường.
    \item \textbf{Độ tin cậy (Resiliency):} Nhờ có đường link dự phòng HA và OSPF, hệ thống đạt được thời gian Uptime cực cao. Việc mất một đường link không còn là thảm họa làm ngưng trệ toàn bộ hoạt động của doanh nghiệp.
    \item \textbf{Bảo mật chiều sâu (Defense-in-Depth):} Bằng cách kết hợp ACL tại Firewall và phân chia mạng thành các VLAN/Subnet riêng biệt, hacker nếu chiếm được một máy trạm Inside cũng khó lòng "pivot" sang các vùng nhạy cảm khác hoặc tấn công vào DMZ nhờ các bộ lọc đã được thiết lập chặt chẽ.
\end{itemize}

\subsection{Kết luận và Hướng phát triển}
Đồ án đã mô phỏng thành công một hệ thống mạng Campus 3 lớp tiêu chuẩn, đáp ứng đầy đủ các tiêu chí từ cơ bản đến chuyên sâu (Mức 3). Hệ thống không chỉ dừng lại ở việc kết nối thông thường mà còn tích hợp các công nghệ hiện đại như: Định tuyến động OSPF với khả năng HA, Bảo mật biên bằng Firewall/NAT/ACL, và đặc biệt là cơ chế tự động hóa mạng (Network Automation) thông qua Script Python để cân bằng tải và giám sát an ninh trực quan (Heatmap).

\textbf{Hướng phát triển tương lai:}
\begin{itemize}
    \item Triển khai thuật toán cân bằng tải phức tạp hơn (như Round-Robin hoặc Least-Connection) trên nền tảng Nginx/HAProxy thay vì dùng Iptables ở lớp mạng.
    \item Áp dụng SDN (Software Defined Networking) với giao thức OpenFlow để điều khiển tập trung luồng dữ liệu thay vì cấu hình tĩnh lẻ tẻ trên từng Switch.
    \item Tích hợp hệ thống phát hiện xâm nhập (IDS/IPS) như Snort để tự động cập nhật các rule ACL khi phát hiện dấu hiệu tấn công bất thường.
\end{itemize}
